
\medskip
\phantomsection\label{sec:sga3-VIA-2.6}%
\noindent\SGAInlineRefTarget{sga3:vi:paragraph:2-6}\textbf{2.6. Complements.---}
This paragraph has been added; it is taken from [Per75], II,
\S\S~1--2, with complements due to O.~Gabber.%
\footnote{\label{note:sga3-VIA-31}Editor's note: these results were
communicated to us by O.~Gabber, in particular \SGARefLink{prop:sga3-VIA-2.6.6}{2.6.6}, which plays an
important role in \hyperref[sec:sga3-VIB-5]{Section~5} of Expos\'e~VI\textsubscript{B}.}
It shows that the preceding results are valid for every group scheme
\(G\) over a field \(k\). (This will be used in \hyperref[sec:sga3-VIB-5]{Sections~5}, 6, and 7
of Expos\'e~VI\textsubscript{B}.)

Fix a field \(k\). We begin with the following lemma
(loc.\ cit., II, 2.1.1), which does not appear explicitly in
EGA~IV\textsubscript{2}, \S~4.4 (although it can probably be read
between the lines at the beginning of loc.\ cit., \S~4.4.1).

\medskip
\phantomsection\label{lem:sga3-VIA-2.6.0}%
\noindent\textbf{Lemma 2.6.0.}
\textit{Let \(X\) be an irreducible \(k\)-scheme, let \(K\) be an
extension of \(k\), and let \(X'\) be an irreducible component of
\(X_K\). The projection \(X'\to X\) is surjective.}

\noindent\textit{Proof.}
Let \(B\) be a transcendence basis of \(K\) over \(k\), and set
\(L=k(B)\subset K\). By EGA~IV\textsubscript{2}, 4.3.2 and 4.4.1,
\(X_L\) is irreducible and \(X'\) dominates \(X_L\). Hence
\(X'\to X_L\) is integral and dominant, and therefore surjective.
Since \(X_L\to X\) is surjective (loc.\ cit., 4.4.1), so is
\(X'\to X\).

For the rest of \hyperref[sga3:vi:paragraph:2-6]{2.6}, fix a \(k\)-group scheme \(G\) and an action
\(\mu:G\times X\to X\) of \(G\) on a \(k\)-scheme \(X\), satisfying
the following condition:
\[
  \tag{\(\star\)}\label{sga3:vi:lemma:2-6-0:equation:star}
  \Phi:G\times X\longrightarrow X\times X,\qquad
  (g,x)\longmapsto(gx,x)
  \quad\text{is surjective.}
\]
(This holds in particular when \(G\) acts on itself by left
translations.) For short, we shall say: ``Let \(X\) be a
\(G\)-scheme satisfying \((\star)\).'' Finally, write \(k'\) for the
perfect closure of \(k\).

\medskip
\phantomsection\label{prop:sga3-VIA-2.6.1}%
\noindent\textbf{Proposition 2.6.1.}
\textit{Let \(X\) be a \(G\)-scheme satisfying \((\star)\).}
\begin{enumerate}
\item[\SGAInlineRefTarget{prop:sga3-VIA-2.6.1:item:i}(i)] \textit{\(X\) is geometrically pointwise irreducible over
\(k\): for every extension \(K\) of \(k\), each point \(x\in X_K\)
belongs to a unique irreducible component of \(X_K\).}
\item[\SGAInlineRefTarget{prop:sga3-VIA-2.6.1:item:ii}(ii)] \textit{Every local ring of \((X_{k'})_{\mathrm{red}}\)
is normal.}
\item[\SGAInlineRefTarget{prop:sga3-VIA-2.6.1:item:iii}(iii)] \textit{Let \(\eta\) be a maximal point of
\((X_{k'})_{\mathrm{red}}\), put
\(C=\overline{\{\eta\}}\), and let \(L\) be the algebraic closure of
\(k\) in \(\kappa(\eta)\). Then \(C\) is an \(L\)-scheme and is
geometrically irreducible over \(L\); that is,
\(C\otimes_LK\) is irreducible for every extension \(K\) of \(L\).}
\item[(iv)] \textit{In particular, if \(x\) is a rational point of
\(X\), then the irreducible component of \(X\) containing \(x\) is
geometrically irreducible over \(k\).}
\end{enumerate}

\noindent\textit{Proof.}
\SGARefLink{prop:sga3-VIA-2.6.1:item:i}{(i)} Since hypothesis \((\star)\) is preserved by every base change
\(k\to K\), it is enough to show that every \(x\in X\) belongs to a
unique irreducible component of \(X\). Since
\(\operatorname{Spec}(k')\to\operatorname{Spec}(k)\) is a universal
homeomorphism, we may further suppose that \(k\) is perfect. We may
then suppose that \(G\) and \(X\) are reduced.

Let \(\eta\) be a maximal point of \(X\), and let \(z\) be any point
of \(Z=\overline{\{\eta\}}\). Since \(X\) is reduced,
\(\mathcal O_{X,\eta}=\kappa(\eta)\); and, since \(k\) is perfect,
\(\kappa(\eta)\otimes_kK\) is normal for every extension \(K\) of
\(k\) (cf.\ EGA~IV\textsubscript{2}, 6.14.2). Thus every point of
\(X_K\) lying above \(\eta\) is normal.

Since \(\Phi\) is surjective, there is a point \(\gamma\) of
\(G\times X\) such that the two projections of \(\Phi(\gamma)\) are
\(z\) and \(\eta\). Set \(K=\kappa(\gamma)\). There are then rational
points \(g\in G_K\) and \(\eta'\in X_K\), with \(\eta'\) lying above
\(\eta\), such that \(z'=g\eta'\) lies above \(z\). By the preceding
paragraph, \(\eta'\) is a normal point of \(X_K\), and therefore so is
\(z'\). Since \(\pi:X_K\to X\) is flat, it follows that
\(z=\pi(z')\) is a normal point of \(X\)
(cf.\ EGA~IV\textsubscript{2}, 2.1.13). This proves (ii) and (i).

The first assertion of \SGARefLink{prop:sga3-VIA-2.6.1:item:iii}{(iii)} now follows from \SGARefLink{prop:sga3-VIA-2.6.1:item:ii}{(ii)}. Since \(L\) is
algebraically closed in \(\kappa(\eta)\), \(C\) is geometrically
irreducible over \(L\), by EGA~IV\textsubscript{2}, 4.5.9.

Finally, if \(X\) has a rational point \(x\), it follows from \SGARefLink{prop:sga3-VIA-2.6.1:item:iii}{(iii)}
that \(L=k\), and hence the irreducible component \(C\) containing
\(x\) is geometrically irreducible over \(k\).\begingroup
\renewcommand{\thefootnote}{\fnsymbol{footnote}}%
\footnote[1]{Translator's note: the French source prints \(X\) in
this conclusion and then switches from the point \(x\) to undefined
symbols \(e\) and \(e_K\). The statement and argument require the
component \(C\) and the points \(x\) and \(x_K\), respectively.}%
\endgroup
This can also be seen directly as follows (cf.\ [Per75], II, \SGARefLink{par:sga3-VIA-2.1}{2.1}).
Let \(C\) be the irreducible component of \(X\) containing \(x\), and
let \(K\) be an extension of \(k\). The scheme \(X_K\) has a unique
point \(x_K\) above \(x\), and by \SGARefLink{prop:sga3-VIA-2.6.1:item:i}{(i)}, \(x_K\) belongs to a unique
irreducible component \(C'\) of \(X_K\). On the other hand, by
\hyperref[lem:sga3-VIA-2.6.0]{2.6.0}, every irreducible component of
\(C_K\) contains \(x_K\), and therefore equals \(C'\).

\medskip
\phantomsection\label{not:sga3-VIA-C0}%
\noindent\textbf{Notation.---}
Write provisionally \(C^0\) for the reduced closed subscheme of \(G\)
whose underlying space is the unique irreducible component of \(G\)
containing the identity element \(e\).

\medskip
\phantomsection\label{cor:sga3-VIA-2.6.2}%
\noindent\textbf{Corollary 2.6.2.}
\textit{\(C^0\) is geometrically irreducible over \(k\) and is
set-theoretically stable under the group law; that is,
\((C^0_{k'})_{\mathrm{red}}\) is a subgroup of
\((G_{k'})_{\mathrm{red}}\). Consequently, \(C^0\) is
quasi-compact.}

\noindent\textit{Proof.}
By \hyperref[prop:sga3-VIA-2.6.1]{2.6.1}, \(C^0\) is geometrically
irreducible over \(k\), so \(C^0\times C^0\) is irreducible and
\[
  \nu(C^0\times C^0)\subset C^0,
\]
where \(\nu\) denotes the morphism \((g,h)\mapsto gh^{-1}\). Since
\(\operatorname{Spec}k'\to\operatorname{Spec}k\) is a universal
homeomorphism, the same conclusion holds for \(C^0_{k'}\), and then
for
\[
  H=(C^0_{k'})_{\mathrm{red}}.
\]
Since \(H\times_{k'}H\) is reduced, \(\nu\) induces a morphism
\[
  H\times_{k'}H\longrightarrow H;
\]
thus \(H\) is a subgroup of \((G_{k'})_{\mathrm{red}}\).
Consequently, by \hyperref[sga3:VIA:0.5.1]{0.5.1}, \(H\), and
hence \(C^0\), is quasi-compact.

\medskip
\phantomsection\label{rem:sga3-VIA-2.6.3}%
\noindent\textbf{Recall 2.6.3.---}
Let \(Y\) be a scheme. Recall (EGA~0\textsubscript{III}, 9.1.1) that
a subset \(E\) of \(Y\) is called \emph{retrocompact} if the
inclusion \(E\hookrightarrow Y\) is quasi-compact. Moreover, by
EGA~IV\textsubscript{1}, 1.9.5(v) and 1.10.1, if \(U\) is a
retrocompact open subset of \(Y\), its closure \(\overline U\) is the
union of the closures \(\overline{\{y\}}\) of the points \(y\in U\).
Of course, it is enough to let \(y\) run through the maximal points of
\(U\), that is, the maximal points of \(Y\) contained in \(U\).

Consequently, if \(Y\) is quasi-separated and \(\eta\) is a maximal
point of \(Y\), then the intersection of the closed neighborhoods of
\(\eta\) equals \(\overline{\{\eta\}}\). Indeed, if
\(y\in Y-\overline{\{\eta\}}\), then \(y\) belongs to an affine open
\(V\) not containing \(\eta\). Since \(Y\) is quasi-separated, \(V\)
is retrocompact, and therefore \(\overline V\) is the union of the
\(\overline{\{\xi\}}\), for \(\xi\) running through the maximal points
of \(Y\) belonging to \(V\). Hence \(\eta\notin\overline V\), so
\(Y-\overline V\) is a closed neighborhood of \(\eta\) not containing
\(y\).

\medskip
\phantomsection\label{prop:sga3-VIA-2.6.4}%
\noindent\textbf{Proposition 2.6.4.}
\textit{Let \(X\) be a \(G\)-scheme satisfying \((\star)\), and let
\(U\) be an open subset of \(X\).}
\begin{enumerate}
\item[\SGAInlineRefTarget{prop:sga3-VIA-2.6.4:item:i}(i)] \textit{\(C^0U\) is an open subset of \(X\), equal to the
union of the irreducible components of \(X\) whose generic point
belongs to \(U\).}
\item[(i\(^{\prime}\))] \textit{Consequently, if \(X\) is
irreducible, it is quasi-compact.}
\item[\SGAInlineRefTarget{prop:sga3-VIA-2.6.4:item:ii}(ii)] \textit{If, moreover, \(U\) is retrocompact in \(X\), then
\(C^0U=\overline U\), and this subset is both open and closed in
\(X\).}
\end{enumerate}

\noindent\textit{Proof.}
\SGARefLink{prop:sga3-VIA-2.6.4:item:i}{(i)} First, \(C^0U\) is open, since it is the union, for \(g\in C^0\),
of the projections of the open subsets
\[
  g\cdot U_{\kappa(g)}\subset X_{\kappa(g)},
\]
and every projection \(X_{\kappa(g)}\to X\) is open.

To prove the second assertion of \SGARefLink{prop:sga3-VIA-2.6.4:item:i}{(i)}, we may replace \(k\) by \(k'\)
(since \(\operatorname{Spec}k'\to\operatorname{Spec}k\) is a
universal homeomorphism), and therefore suppose that \(k\) is perfect.
We may then suppose that \(G\) and \(X\) are reduced, and hence
geometrically reduced.

Let \(\eta\) be a maximal point of \(X\) contained in \(U\), and put
\(Z=\overline{\{\eta\}}\). Consider
\[
  \mu':C^0\times Z\longrightarrow X.
\]
Since \(C^0\) is geometrically irreducible, \(C^0\times Z\) is
irreducible; write \(\gamma\) for its generic point. The morphism
\(\mu'\) sends the point \(\varepsilon(\eta)\), where \(\varepsilon\)
denotes the identity section of \(C^0\), to \(\eta\). Hence
\(\gamma\) is sent to a generization of \(\eta\), and therefore to
\(\eta\). Thus \(\mu'\) sends the underlying space of \(C^0\times Z\)
into \(Z\), and, since \(C^0\times Z\) is reduced, \(\mu'\) factors
through \(Z\).

Now let \(z\in Z\), and put \(K=\kappa(z)\). The morphism
\[
  \mu_z:G_K\longrightarrow X_K,\qquad h\longmapsto h\cdot z,
\]
is surjective. Let \(\alpha\) be a maximal point of \(X_K\). The local
ring \(\mathcal O_{X_K,\alpha}\) is a field, since \(X_K\) is reduced,
so \(\mu_z\) is flat at every point of \(G_K\) above \(\alpha\).
It follows from \hyperref[lem:sga3-VIA-2.5.3]{Lemma~2.5.3} that
\(\mu_z\) is flat.

On the other hand, \(\mu_z\) sends the generic point \(\omega\) of
\(C^0_K\) to a point \(t\in Z_K\). Let \(\beta\) be a maximal point
of \(Z_K\) such that \(t\in\overline{\{\beta\}}\). Since \(\mu_z\) is
flat, there is a generization \(\xi\) of \(\omega\) such that
\(\mu_z(\xi)=\beta\). But \(\omega\) is a maximal point of \(G_K\), so
necessarily \(\xi=\omega\). Thus \(\mu_z(\omega)=\beta\), and
\(\beta\) lies above \(\eta\), because \(Z_K\to Z\) is flat.

Put \(L=\kappa(\omega)\), and let \(\omega_L\) and \(z_L\) be the
\(L\)-points obtained from \(\omega\) and \(z\). Then
\[
  \omega_L\cdot z_L=\beta'
\]
is a point of \(Z_L\) lying above \(\beta\), and hence
\[
  z_L=\omega_L^{-1}\cdot\beta'\in C^0_L\cdot U_L.
\]
Therefore \(z\in C^0\cdot U\), which proves \SGARefLink{prop:sga3-VIA-2.6.4:item:i}{(i)}.

Since \(C^0\) is quasi-compact by
\hyperref[cor:sga3-VIA-2.6.2]{2.6.2}, assertion
(i\(^{\prime}\)) follows: if \(X\) is irreducible and \(U\) is a
nonempty affine open subset, then \(X=C^0U\), so \(X\) is the image of
the morphism \(C^0\times U\to X\), and is therefore quasi-compact.
Finally, if \(U\) is retrocompact in \(X\), then by
\hyperref[rem:sga3-VIA-2.6.3]{2.6.3}, \(\overline U\) is the union of
the \(\overline{\{\eta\}}\), where \(\eta\) runs through the maximal
points of \(X\) contained in \(U\). Hence \(\overline U=C^0U\).
This proves \SGARefLink{prop:sga3-VIA-2.6.4:item:ii}{(ii)}.
